\documentclass[11pt,a4paper]{article}

\usepackage[T1]{fontenc}
\usepackage[utf8]{inputenc}
\usepackage{lmodern}
\usepackage{microtype}
\usepackage[margin=1.02in]{geometry}
\usepackage{amsmath,amssymb,amsthm,mathtools}
\usepackage[hidelinks]{hyperref}
\hypersetup{
  pdftitle={A Universal Square-Root-of-Six Bound for Kissing Numbers},
  pdfauthor={Sebastien Palcoux}
}
\usepackage{xurl}

\newtheorem{theorem}{Theorem}
\newtheorem{lemma}[theorem]{Lemma}
\newtheorem{corollary}[theorem]{Corollary}
\theoremstyle{remark}
\newtheorem{remark}[theorem]{Remark}

\newcommand{\R}{\mathbb{R}}
\newcommand{\vol}{\operatorname{vol}}

\title{A Universal Square-Root-of-Six Bound for Kissing Numbers}
\author{S\'ebastien Palcoux}
\date{September 6, 2026}

\begin{document}
\maketitle

\begin{abstract}
Let $\tau_n$ denote the kissing number in Euclidean dimension $n$. We give two elementary proofs that
\[
\tau_n\leq (\sqrt6)^n\qquad(n\geq1).
\]
The first uses spherical-cap areas and a short two-step recurrence. The second packs explicit bodies of revolution inside the unit ball and is the route used by an end-to-end Lean~4/Mathlib certificate registered in Palomar. Since $\tau_2=6$, the base $\sqrt6$ is optimal for a bound of this form valid in every dimension. The cap proof is not part of the formal certificate.
\end{abstract}

\begin{center}
\begin{minipage}{0.88\textwidth}
\small
\textbf{AI provenance and verification status.}
The mathematical arguments, manuscript, and Lean~4/Mathlib formalization were generated with OpenAI's GPT-6 Astra, from prompts supplied by the author. The volume-based formalization was registered in Palomar on September~5, 2026, as \href{https://palomar-registry.org/entry.html?id=PALOMAR-2026-09-05-000002&version=1}{\nolinkurl{PALOMAR-2026-09-05-000002}, version~1}, following successful mechanical verification and an automated review that identified no problems. This later revision adds the unformalized cap proof and simplifies the exposition; the Lean sources are unchanged. Registration is not human peer review or a novelty certificate. The manuscript has not undergone independent human mathematical review.
\end{minipage}
\end{center}

\section{Statement and a common geometric observation}

A \emph{kissing configuration} in $\R^n$ is a finite set $X$ satisfying
\begin{equation}\label{eq:configuration}
 \|x\|=1\quad(x\in X),\qquad
 \langle x,z\rangle\leq\frac12\quad(x,z\in X,\ x\ne z).
\end{equation}
Multiplying its vectors by two gives the centers of nonoverlapping unit balls touching the unit ball at the origin. We write $\tau_n$ for the largest cardinality of a kissing configuration in $\R^n$. The bounds below also establish the existence of this maximum.

\begin{theorem}\label{thm:main}
For every integer $n\geq1$,
\[
\tau_n\leq(\sqrt6)^n.
\]
\end{theorem}

\begin{corollary}\label{cor:base}
The least real number $\alpha>0$ such that $\tau_n\leq\alpha^n$ for every $n\geq1$ is $\sqrt6$. Equivalently,
\[
\sup_{n\geq1}\tau_n^{1/n}=\sqrt6.
\]
\end{corollary}

The point is the optimal base in a bound valid in every dimension with no additional prefactor. No claim of asymptotic optimality is intended.

Throughout, put $s=\sqrt3$ and $a=s/2$. For a unit vector $x$, let
\[
 C_x=\{y\in\R^n:\langle x,y\rangle>a\|y\|\}.
\]
This is the open cone of half-angle $\pi/6$ about $x$; its intersection with the unit sphere is the corresponding open spherical cap.

\begin{lemma}\label{lem:separation}
For distinct members $x,z$ of a kissing configuration, $C_x\cap C_z=\varnothing$.
\end{lemma}

\begin{proof}
A vector in the intersection could be normalized to a unit vector $u$ satisfying
\[
 s<\langle x+z,u\rangle\leq\|x+z\|
 =\sqrt{2+2\langle x,z\rangle}\leq\sqrt3=s,
\]
a contradiction.
\end{proof}

Both proofs use this observation, but otherwise establish the bound independently. In dimension one, every unit vector is $1$ or $-1$, so $|X|\leq2<\sqrt6$; henceforth we consider $n\geq2$.

\section{First proof: a spherical-cap recurrence}\label{sec:cap}

Let $c_n$ be the fraction of the area of $S^{n-1}$ occupied by a cap of angular radius $\pi/6$. Polar coordinates about its center give
\begin{equation}\label{eq:capfraction}
 c_n=\frac{I_{n-2}}{J_{n-2}},\qquad
 I_m=\int_0^{\pi/6}\sin^m t\,dt,\qquad
 J_m=\int_0^\pi\sin^m t\,dt.
\end{equation}
Indeed, the surface element is a constant multiple of $\sin^{n-2}t\,dt$, and the constant cancels. We use open caps, so Lemma~\ref{lem:separation} gives, for every kissing configuration $X\subset\R^n$,
\begin{equation}\label{eq:cappacking}
 |X|c_n\leq1.
\end{equation}

\begin{lemma}\label{lem:caprecurrence}
For every integer $n\geq2$,
\begin{equation}\label{eq:caprecurrence}
 c_{n+2}\geq\frac{n}{4(n+1)}c_n\geq\frac{c_n}{6}.
\end{equation}
\end{lemma}

\begin{proof}
Write $m=n-2$. Integration by parts gives
$J_{m+2}=(m+1)J_m/(m+2)$. On substituting $u=\sin t$, we have
\[
 I_m=\int_0^{1/2}\frac{u^m}{\sqrt{1-u^2}}\,du.
\]
A second integration by parts, with the boundary term
$u^{m+1}(1-4u^2)/\sqrt{1-u^2}$ vanishing at both endpoints, yields
\[
 4(m+3)I_{m+2}-(m+1)I_m
 =\int_0^{1/2}\frac{u^{m+2}(1-4u^2)}{(1-u^2)^{3/2}}\,du\geq0.
\]
Consequently
\[
 \frac{c_{n+2}}{c_n}
 =\frac{m+2}{m+1}\frac{I_{m+2}}{I_m}
 \geq\frac{m+2}{4(m+3)}
 =\frac{n}{4(n+1)}\geq\frac16.
\]
\end{proof}

\begin{proof}[First proof of Theorem~\ref{thm:main}]
Direct integration gives
\[
 c_2=\frac16,\qquad c_3=\frac{2-s}{4}>\frac1{15},
\]
where $s<26/15$ follows by squaring. The recurrence therefore gives
\[
 c_n\geq6^{-n/2}\qquad(n\geq2\text{ even}).
\]
For the odd dimensions, its stronger form at $n=3$ gives
\[
 c_5\geq\frac3{16}c_3>\frac1{80}>6^{-5/2},
\]
since $80^2<36^2\cdot6$. Further applications of the recurrence give
$c_n>6^{-n/2}$ for every odd $n\geq5$. Equation~\eqref{eq:cappacking} proves the required bound in these dimensions. In dimension three it gives $|X|<15$, hence
\[
 |X|\leq14<6\sqrt6=(\sqrt6)^3.
\]
Together with the dimension-one observation, this bounds every configuration. Its realizable cardinalities form a nonempty finite set of nonnegative integers, so their maximum $\tau_n$ exists and satisfies the same bound.
\end{proof}

\section{Second proof: explicit Euclidean volumes}\label{sec:volume}

This is a condensed presentation of the construction used in the registered Lean certificate. It does not use the cap-area formula or the recurrence of Section~\ref{sec:cap}.

Let $B_n$ be the open unit ball and $V_n=\vol(B_n)$. Fix $e=(1,0)$ in orthogonal coordinates $(t,v)\in\R\times\R^{n-1}$. Rotating any measurable body $K\subset C_e\cap B_n$ toward the points of $X$ gives disjoint bodies of equal volume by Lemma~\ref{lem:separation}. Thus
\begin{equation}\label{eq:volpacking}
 |X|\vol(K)\leq V_n.
\end{equation}
For a continuous nonnegative profile $g$ on $[b,c]$, Fubini's theorem gives
\begin{equation}\label{eq:profilevolume}
 \vol\{(t,v):b<t<c,\ \|v\|<g(t)\}
 =V_{n-1}\int_b^c g(t)^{n-1}\,dt.
\end{equation}
The coordinate conditions for containment are
\begin{equation}\label{eq:coordinates}
 (t,v)\in C_e\iff t>0\text{ and }3\|v\|^2<t^2,
 \qquad (t,v)\in B_n\iff t^2+\|v\|^2<1.
\end{equation}

\subsection{A bicone and the dimensions \texorpdfstring{$n\geq5$}{n >= 5}}

Consider the union $K_n=L_n\cup U_n$, where
\begin{align*}
 L_n&=\{(t,v):0<t<a,\ \|v\|<t/s\},\\
 U_n&=\{(t,v):a<t<1,\ \|v\|<(2+s)(1-t)\}.
\end{align*}
It is a bicone of total height one and common base radius $1/2$, with the base section omitted. More precisely, each point is a strict convex combination of a point of the disk
\[
 D=\{(a,v):\|v\|<1/2\}
\]
with either $0$ or $e$. The disk lies in $C_e\cap B_n$, whereas $0,e$ satisfy the weak inequalities defining their closures. The triangle inequality therefore gives $K_n\subset C_e\cap B_n$.

By \eqref{eq:profilevolume}, the two cones have volumes
$aV_{n-1}/(n2^{n-1})$ and $(1-a)V_{n-1}/(n2^{n-1})$. Hence
\begin{equation}\label{eq:bicone}
 \vol(K_n)=\frac{V_{n-1}}{n2^{n-1}},\qquad
 |X|\leq A_n:=n2^{n-1}\frac{V_n}{V_{n-1}}.
\end{equation}
The standard ball-volume formula gives
\begin{equation}\label{eq:ballrecurrence}
 V_d=\frac{\pi^{d/2}}{\Gamma(d/2+1)},\qquad
 V_{d+2}=\frac{2\pi}{d+2}V_d,
\end{equation}
and consequently
\begin{equation}\label{eq:Anrecurrence}
 A_{n+2}=\frac{4(n+1)}{n}A_n\leq6A_n\qquad(n\geq2).
\end{equation}
The initial estimates
\[
 A_5=\frac{256}{3}<36\sqrt6=(\sqrt6)^5,\qquad
 A_6=60\pi<216=(\sqrt6)^6
\]
follow by squaring the first comparison and using $\pi<22/7$ for the second. Induction on each parity proves $|X|<(\sqrt6)^n$ for all $n\geq5$. In dimension two, the same bicone gives $|X|\leq A_2=2\pi<7$, so integrality gives $|X|\leq6$.

\subsection{Dimension three}

Retain $L_3$ and replace the upper cone by
\[
 W_3=\{(t,v):a<t<1,\ \|v\|<\sqrt{1-t^2}\}.
\]
It lies in the ball by construction, and in the cone because $t>a$ implies $3(1-t^2)<t^2$. Since $V_2=\pi$, formula~\eqref{eq:profilevolume} gives
\begin{align*}
 \vol(L_3\cup W_3)
 &=\pi\left(\frac{s}{24}+\int_a^1(1-t^2)\,dt\right)
 =\frac{(2-s)\pi}{3}
 >\frac{4\pi}{45}=\frac{V_3}{15}.
\end{align*}
Here $V_3=4\pi/3$, and the strict inequality uses $s<26/15$. Thus \eqref{eq:volpacking} yields $|X|<15$, hence $|X|\leq14<6\sqrt6$.

\subsection{Dimension four}

A polynomial profile avoids a trigonometric volume integral. Retain $L_4$ and put
\[
 p(t)=(1-t)(14t+2-6s),\qquad
 W_4=\{(t,v):a<t<1,\ \|v\|<p(t)\}.
\]
For $a\leq t\leq1$ we have $p(t)\geq0$, since $14t+2-6s\geq s+2>0$. Writing $u=t-a$, a factorization gives
\begin{equation}\label{eq:fourcontainment}
 1-t^2-p(t)^2
 =(1-t)(t-a)\bigl(4s-6+(126s-140)u+196u^2\bigr)\geq0.
\end{equation}
Indeed, $u\geq0$ and the constant and linear coefficients in the last factor are positive because $s>3/2$. This puts $W_4$ in the ball; as in dimension three, $t>a$ then puts it in the cone.

For completeness, its volume requires only the following polynomial integral. Set $\delta=1-a$ and $b=16-6s$. Substituting $u=1-t$ gives
\begin{align*}
 \int_a^1p(t)^3\,dt
 &=\int_0^\delta u^3(b-14u)^3\,du\\
 &=\frac{b^3\delta^4}{4}-\frac{42b^2\delta^5}{5}
       +98b\delta^6-392\delta^7
 =\frac{33169-19150s}{40}.
\end{align*}
Adding the lower cone, whose profile integral is $s/64$, yields
\begin{equation}\label{eq:fourvolume}
 \vol(L_4\cup W_4)=T V_3,\qquad
 T=\frac{265352-153195s}{320}.
\end{equation}
The rational bound $s<1351/780$ follows from $1351^2-3\cdot780^2=1$. Therefore
\[
 T>\frac{541}{16640}>\frac{33}{1036}
 =\frac{3}{296}\frac{22}{7}>\frac{3\pi}{296}.
\]
Since $V_4=\pi^2/2$, this proves
\[
 \vol(L_4\cup W_4)>\frac{3\pi}{296}\frac{4\pi}{3}
 =\frac{\pi^2}{74}=\frac{V_4}{37}.
\]
Equation~\eqref{eq:volpacking} now gives $|X|<37$, hence $|X|\leq36=(\sqrt6)^4$.

\begin{proof}[Completion of the second proof of Theorem~\ref{thm:main}]
The volume estimates cover every $n\geq2$, and the dimension-one observation covers $n=1$. Thus every kissing configuration satisfies the claimed bound, independently of the first proof. As before, the finite nonempty set of realizable cardinalities has an attained maximum, which is $\tau_n$.
\end{proof}

\section{Optimality of the universal base}

\begin{proof}[Proof of Corollary~\ref{cor:base}]
The six points
\[
 (1,0),\quad (\tfrac12,\tfrac{s}{2}),\quad
 (-\tfrac12,\tfrac{s}{2}),\quad (-1,0),\quad
 (-\tfrac12,-\tfrac{s}{2}),\quad (\tfrac12,-\tfrac{s}{2})
\]
are distinct unit vectors whose pairwise inner products are at most $1/2$. They form a kissing configuration, and Theorem~\ref{thm:main} gives the reverse bound, so $\tau_2=6$.

The base $\sqrt6$ works by the theorem. Conversely, any positive universal base $\alpha$ must satisfy $6=\tau_2\leq\alpha^2$, hence $\alpha\geq\sqrt6$. Likewise, $\tau_n^{1/n}\leq\sqrt6$ in every dimension, with equality at $n=2$, proving the supremum statement.
\end{proof}

\section{Formalization status and Palomar registration}\label{sec:lean}

\paragraph{Why the shorter proof is not yet formalized.}
The cap proof uses a familiar formula for spherical area and a few integral identities. To check this route, these steps must be connected to the precise geometric definitions and supplied as verified lemmas; that work has not been completed in this project. This is unfinished formalization work, not an obstruction in Lean. The volume proof avoids that dependency and already gives a complete certificate for the theorem and its optimality.

\paragraph{What the certificate establishes.}
The accompanying repository is
\begin{center}
\url{https://github.com/sebastienpalcoux/sqrt6-kissing-bound}.
\end{center}
Its certificate follows the Euclidean-volume route of Section~\ref{sec:volume}: cone separation, measurable disjoint packing, the explicit body-volume integrals, and every dimension case. It also proves that the kissing number is an attained maximum and verifies the six-point configuration used for optimality. Standard analysis comes from proved Mathlib theorems. The present exposition shortens the separation argument and combines the high-dimensional estimates into \eqref{eq:Anrecurrence}; the Lean sources retain their original intermediate lemmas. The certificate does not check the cap-area formula or Lemma~\ref{lem:caprecurrence}.

The eight principal declarations selected by \texttt{comparator.json} were registered as
\begin{center}
\href{https://palomar-registry.org/entry.html?id=PALOMAR-2026-09-05-000002&version=1}{\texttt{PALOMAR-2026-09-05-000002}, version~1}.
\end{center}
They express the universal bound, attained kissing numbers, the planar value six, and the equivalent optimality statements. The registered repository commit is
\begin{center}
\small\texttt{cc87870ad897f0fe15e3cdffd3d7acd15f0c1ba9}.
\end{center}
Comparator checked the correspondence between the independently stated Challenge and its Solution. Lean's kernel and the independent NanoDa checker verified the exported proofs, and Palomar's automated editorial review identified no problems. These checks do not constitute independent human review or establish novelty.

The registered development pins Lean~4.33.0 and Mathlib revision
\begin{center}
\small\texttt{db584cd6d46c92f209a44c0f1c829460d327499d}.
\end{center}
Palomar preserves that exact source snapshot, including its original manuscript. This later two-proof revision does not change the registered version, Lean sources, Comparator interface, or dependency pins.

After obtaining the dependencies with \texttt{lake exe cache get}, run
\begin{center}
\begin{tabular}{l}
\texttt{bash scripts/check.sh}\\
\texttt{bash scripts/verify-comparator.sh}
\end{tabular}
\end{center}
These commands build the certificate and audit its axioms, then repeat the Comparator and NanoDa checks. Only Lean's standard foundational axioms \texttt{propext}, \texttt{Classical.choice}, and \texttt{Quot.sound} are allowed. The isolated Challenge contains eight deliberate statement placeholders, whose replacements are verified by Comparator; the Solution and substantive proofs contain no admitted results or additional axioms.

\section*{Acknowledgment}
This note was prompted by the MathOverflow question ``Upper bound of the kissing number in $n$ dimensions'' and by Henry Cohn's answer explaining how effective spherical-code bounds bear on the problem.

\begingroup\small
\begin{thebibliography}{9}
\bibitem{MO-question}
S.~Palcoux,
\emph{Upper bound of the kissing number in $n$ dimensions},
MathOverflow question 282644 (2017),
\url{https://mathoverflow.net/questions/282644/}.

\bibitem{Cohn-answer}
H.~Cohn,
\emph{Answer to ``Upper bound of the kissing number in $n$ dimensions''},
MathOverflow answer 282685 (2017),
\url{https://mathoverflow.net/a/282685/}.
\end{thebibliography}
\endgroup

\vspace{0.5em}
\noindent
\textsc{Beijing Institute of Mathematical Sciences and Applications (BIMSA),}\newline
\textsc{Huairou District, Beijing 101408, China}\newline
\texttt{sebastienpalcoux@gmail.com}

\end{document}
